\newpage
\section{Laboratory Work and Experimental Results}

\subsection{Initial Hardware Testing}

\subsubsection{LQE Validation}

\par Before testing the performance of the LQG controller, the performance of the LQE was strictly 
validated on the physical hardware. The primary objective was to ensure that the estimated state 
vector, $\hat{x}$, accurately tracked the physical system's dynamics—specifically the angular 
positions $\alpha$ and $\beta$

\par A fundamental observation regarding the experimental setup was the high precision and 
discrete nature of the sensors. The Quanser platform utilizes a potentiometer for the rotary arm 
and a high-resolution optical encoder for the pendulum angle. Due to their discrete characteristics 
and high accuracy, the measurement noise was found to be negligible, effectively implying a standard 
deviation of zero in ideal conditions. However, to maintain numerical robustness in the Kalman Filter 
algorithm and reflect a high degree of confidence in the hardware readings, the measurement noise 
covariance matrix was set to:

\begin{equation}
    R_e = I_2 \cdot (\text{deg2rad}(0.1))^2
\end{equation}

\par This value corresponds to a very small standard deviation of 0.1°, prioritizing the sensor 
data over the model's predictions. To optimize the estimator's transient response, the process noise 
covariance matrix, $Q_e$, was tuned through empirical testing on the hardware.

\par Figure [X] presents a comparative analysis of the LQE performance for three different 
weights: $Q_e = 0.1 \cdot I_5$, $Q_e = I_5$ and $Q_e = 10 \cdot I_5$. By examining the plots of the 
measured angles ($\alpha, \beta$) against their respective estimates ($\hat{\alpha}, \hat{\beta}$), 
it is evident that the configuration with $Q_e = 10 \cdot I_5$ provided superior tracking fidelity. 
This setup allowed the estimates to converge rapidly to the physical values with minimal lag, 
ensuring that the velocity states (which are not directly measured) were also accurately derived. 
Based on this visual validation, the more aggressive tuning was selected for the subsequent stages 
of the project.

\subsubsection{LQG Testing}

\subsubsection{Analysis and Problems}


\subsection{Propused Solutions and Simulation}

\subsection{Final Hardware Testing and Validation}
